\section{Discussion}\label{sec:discussion}

\subsection{What the evidence supports}
\emph{The evaluation protocol outweighs the model.} Across our experiments, the choice among LightGBM, XGBoost, Random Forest and their average moves illicit F1 by a few hundredths. The choice between a temporal and a random split moves it by 0.12 for transactions and by 0.3 to 0.6 for wallets (Table~\ref{tab:split}). A reader comparing two papers on these data should therefore first ask how rows were assigned to training and test, and only then which model was used. For wallet-level data in particular, a result is not interpretable unless the paper states what a row is, whether an address can occur on both sides of the split, and whether features are computed per appearance or per address.

\emph{Graph information helps where neighbours carry the label.} The wallet result is the one substantive modelling finding. It required no graph neural network: one step of unweighted neighbour averaging plus a handful of degree statistics, fed to a tree ensemble, accounts for the improvement, and the same construction does almost nothing for transactions. The difference follows from properties of the data that can be checked in advance, namely whether the base features already encode neighbourhoods and how strongly labels cluster along edges (Table~\ref{tab:homophily}). We regard this as a practical suggestion for similar problems: measure label homophily on the training period before investing in graph models. Our finding is compatible with reports that graph networks do not beat feature baselines on Elliptic transactions~\cite{weber2019,maganti2026}; it adds that the conclusion should not be carried over to the wallet level without testing.

\emph{Static detectors have a shelf life, and the remedy has two parts.} The shutdown at step 43 is a single event in a single data set, and we do not generalise its magnitude. Its anatomy is nevertheless instructive. Before the event the model was excellent by every measure, and nothing in its validation performance warned of what followed. Afterwards the failure was complete and was a failure of ranking. Retraining with recent labels repaired the ranking within a few steps but left the fixed decision threshold wrong. An operational system therefore needs both periodic retraining and periodic re-estimation of thresholds, and it needs monitoring that reports a ranking metric and a decision metric side by side; either alone gives a misleading picture (Section~\ref{sec:res-drift}).

\emph{Unsupervised scores need labels to be trusted.} The Isolation Forest result is a caution against treating ``anomaly detection'' as synonymous with ``detection of illicit activity''. In the transaction data the unsupervised score was strongly and stably anti-correlated with the label, a fact that is only discoverable with labels.

\subsection{Unexpected findings}
Four results ran against our expectations and are reported as such. (i) Graph-derived features alone matched the full feature set for wallets, which means that the released per-wallet features add little once neighbourhood structure is known. (ii) The ensemble was not the best model in the network layer, where a default Random Forest out-ranked it. (iii) The three-way verdict collapsed to two classes for transactions, because validation data were too easy to place a review threshold below the decision threshold. (iv) A domain classifier could separate any two periods perfectly, including periods inside the training set, which invalidated it as a drift diagnostic and revealed that the anonymised features encode time.

\subsection{Threats to validity}\label{sec:threats}
\subsubsection{Internal validity} All hyper-parameters were fixed in advance and nothing was tuned on test data, but the supplementary analyses were designed after the main results were known and are exploratory. The validation periods are short (five time steps; one day), which makes threshold selection and early stopping noisy (Section~\ref{sec:res-verdict}); our remedy, averaging over models and seeds, reduces but does not remove this. Bootstrap intervals describe sampling variability of the test set for fixed models and do not include training variability or the choice of split point; we use one split per layer. Entities in the same time step are not independent, so the intervals are likely too narrow.

\subsubsection{Construct validity} Labels in Elliptic++ cover 23\% of transactions and 32\% of wallets and derive from a proprietary process~\cite{weber2019}; metrics describe the labelled subset, and verdicts on unlabelled entities cannot be checked. Wallet features are aggregates over an address's whole observed life. We showed that this affects few test wallets (1.6\% are active in more than one step) and that neighbourhood look-ahead does not drive the results, but the wallet layer remains an evaluation on retrospective features, not a simulation of online scoring. Network labels come from injected attacks and from flags on real traffic~\cite{macia2018}; windows labelled only by blacklist membership were excluded, and the \texttt{anomaly-spam} label is counted as an attack although it does not occur on the test day.

\subsubsection{External validity} Each Bitcoin result rests on one data set covering 49 time steps, and the drift analysis on one event. The network layer rests on five days from an undocumented sample of UGR'16: one test day, 769 positive windows, 87\% of them from one botnet. We cannot say how the sample was drawn, and windows built from sampled flows under-count each host's activity. The network results should be read as a demonstration that the protocol transfers and as a description of this sample, not as a benchmark on UGR'16. The three layers are independent; no claim is made about combining them, and UGR'16 is not Bitcoin network traffic.

\subsubsection{Reproduction of prior work} Our reproduction of the baseline of~\cite{elmougy2023} follows the configuration stated in that paper (50 trees, default settings, steps 1--34 and 35--49) but not its preprocessing in every detail; for example, we do not rescale features, which does not affect tree models. The discrepancy at the wallet level is too large to be explained by such details, but its cause in the original work can only be established by its authors.

\subsection{Limitations of the method and future work}
The prequential analysis assumes that labels arrive immediately, whereas in practice confirmation of illicit activity takes time~\cite{dalpozzolo2018}; repeating it with realistic delays is the most direct extension. Thresholds could be re-estimated on a sliding validation window, and probabilities recalibrated after each retraining~\cite{niculescu2005}. The network configuration should be revisited for very low prevalence and evaluated on complete weeks of UGR'16 with a test week separate from the training week. Finally, the wallet finding invites a comparison with learned graph models under the causal-neighbourhood protocol used here.

\section{Conclusion}\label{sec:conclusion}
We evaluated illicit-activity detectors for Bitcoin transactions, Bitcoin wallets and network flows under one leakage-controlled temporal protocol. The modelling is ordinary; the findings concern what a careful evaluation reveals. On transactions, a tuned ensemble equals but does not exceed a default Random Forest. On wallets, simple graph-derived features double the F1 score, and a prominent published baseline cannot be reproduced under time-ordered splits while a random split approaches it. Over time, a static model fails completely after a regime change, by ranking; retraining restores the ranking and leaves the threshold to be fixed. Unsupervised detection fails in all three layers. Network-flow detection on the available sample is weak and dominated by one attack type. The practical lessons are to state the unit and ordering of a split before any score, to measure label homophily before choosing a graph method, and to monitor ranking and decision quality separately once a detector is deployed.

\section*{Reproducibility}
All experiments are run by scripts in the accompanying repository. \texttt{run\_all.py} executes the six main configurations and writes metrics, tables, figures and per-entity verdicts to \texttt{outputs/}. The scripts in \texttt{paper/analysis/} produce the supplementary analyses and the figures of this paper from those outputs. Random seeds are fixed (base seed 42). The data are public: Elliptic++ at \url{https://github.com/git-disl/EllipticPlusPlus} and the UGR'16 sample at the address given in~\cite{ugrsample}.
